\documentclass[report.tex]{subfiles}

\begin{document}

\subsection*{General Context}

Modern software relies on high-level languages
and formal verification techniques,
such as deductive verification,
model checking,
or static analysis,
to establish correctness guarantees.
Hardware executes machine code produced by compilers,
which translate source programs through intermediate representations
down to assembly.
Consequently,
source-level correctness guarantees remain meaningful only if compilation
strictly preserves program semantics.
Unverified compilation passes risk introducing bugs
or invalidating source-level formal proofs.

Initial efforts in verified compilation established methodologies
for proving semantic preservation.
The CompCert compiler developed by \textcite{leroy-09-compcert}
relies on simulation diagrams,
establishing a correspondence between execution steps
of the source and target programs
while maintaining a state invariant.
CakeML,
a verified compiler for functional languages
made by \textcite{tan-19-cakeml},
combines simulation diagrams with equational reasoning
to prove its correctness.

Beyond general-purpose compilers,
other verified compiler projects focus on specialized application domains.
For instance,
Jinja~\parencite{klein-06-jinja}
targets object-oriented programs compiled to the Java Virtual Machine,
whereas Jasmin~\parencite{almeida-17-jasmin}
verifies a dedicated compiler
for high-performance cryptographic software.

Existing verification paradigms nevertheless exhibit limitations.
On one hand,
equational theories are often incomplete,
especially when programs differ in their memory representations.
On the other hand,
traditional simulation diagrams grow complex
when program control flow structures diverge significantly,
for instance after nested loop exchanges.

Recent advances in formal verification address these shortcomings.
Modern coinductive simulation techniques,
such as those developed by \textcite{cho-23-stuttering}
in the \emph{Stuttering for Free} framework
or implemented by \textcite{xia-20-itrees}
within the \emph{Interaction Trees} library,
simplify the reasoning about program pairs
whose execution steps do not correspond one-to-one.

In parallel,
\emph{relational separation logics}
have emerged as a technique
for local reasoning across dual executions,
allowing memory resources to be tracked independently
or partitioned between two programs.
For instance,
the \emph{Simuliris} framework by \textcite{gaher-22-simuliris}
uses a relational separation logic
to verify source-to-source program transformations.
Recent work on the \emph{Tail Modulo Cons} optimization
by \textcite{allain-25-tmc}
employs relational separation logic
to validate the soundness of the optimization
within an idealized framework.

\subsection*{Research Problem}

This internship investigates whether recent developments
in simulation relations and separation logic
can be integrated into a unified relational separation logic framework.
To our knowledge,
no existing verified compiler incorporates a relational separation logic
to prove the soundness of its optimization passes.
Consequently,
establishing such a framework would enable compositional compiler verification
with improved scalability and fewer language restrictions.
In short,
this approach relaxes constraints present in existing systems,
such as the determinism requirements imposed by CompCert and CakeML.

\subsection*{Contributions}

We developed a logical framework addressing the objectives of this internship.
At its core,
the framework relies on a coinductive simulation relation
capable of matching execution steps of programs with different control flow.
The simulation used is skewed:
it permits exploiting undefined behaviors in the source program
to enable further target-level optimizations.
This simulation serves as the foundation of our verification framework,
defining how source and target execution steps are related.

Building upon this simulation relation,
we defined a relational separation logic
that enables local reasoning over fragments of source and target heaps.
The logic is parameterized by the operational semantics
of the two languages under comparison.
Our framework uses this logical layer to express
and prove relational specifications between program pairs.

We instantiated this generalized framework with a low-level language
that mirrors the RTL language (Register Transfer Language),
commonly used in compiler backends.
We defined the small-step operational semantics for our version of RTL,
with support for arithmetic operations,
dynamic memory allocation and deallocation,
and function calls.

Finally,
we established a system of reasoning rules for RTL programs.
This proof system comprises two distinct categories of rules.
First,
we proved the soundness of asynchronous rules,
which deal with unilateral execution steps taken independently
by either the source or the target.
In particular,
these rules enable target optimizations to exploit undefined behaviors
in the source program.
Second,
we derived synchronous rules,
which deal with joint execution steps when both source and target programs
execute matching instructions at the same time.

\subsection*{Arguments Supporting Its Validity}

We evaluated our framework by proving the correctness
of three program transformations.
The first one is a branch tunneling pass that eliminates redundant control-flow jumps
by collapsing chains of intermediate nodes.
The second one is an instruction-level optimization pass that removes statically known
redundant operations.
Finally, we used our logic to establish the soundness of localized transformations,
such as hoisting a memory read instruction out of a loop body.

Beyond these practical case studies,
we established metatheoretic properties confirming framework soundness.
We proved that the underlying simulation relation is sound,
guaranteeing that it relates only programs exhibiting related behaviors.
Additionally,
we proved our simulation formulation equivalent
to established simulation definitions used in the
literature~\parencites{allain-25-tmc,leroy-09-compcert,gaher-22-simuliris}.

The entire framework has been mechanized
in the Rocq proof assistant~\parencite{rocq}.
All definitions and theorems are machine checked,
relying on the \texttt{coinduction}
Rocq library~\parencite{rocq-coinduction}
and the Iris/stdpp~\parencite{jung-18-iris,rocq-stdpp}
data structures.
A subset of the formalization relies on two standard classical axioms,
specifically the Law of Excluded Middle
and Classical Choice.

\subsection*{Future Work}

This work provides a first step for several research perspectives.
While current case studies verify control flow transformations,
future effort will target data flow optimizations,
such as constant propagation
and common subexpression elimination.
Such passes will test the logical framework when memory states are altered
by eliminated loads, stores or allocations.

A second direction involves supporting multi-language verification.
Our framework is parameterized to relate distinct source and target languages
across translation passes,
but the present formalization operates on a single intermediate language.
Future effort will test the limits of our framework when translating
programs from one language to another.

Finally,
we aim to increase the expressiveness of the underlying languages.
While our logic already accommodates nondeterminism,
it is not present in our current semantics of RTL.
We plan to extend our language semantics with it
to work toward concurrent execution features.
Additionally,
extending the framework from imperative languages to functional paradigms
represents another potential direction for future work.

\end{document}

% LocalWords:  Coinductive coinductive asymmetricity arities
